\chapter{Methods}
\label{ch:methods}

\section{Overview of the pipeline}
\label{sec:methods-overview}

The pipeline follows Ra et al.~\citep{ra2025}, adapted from mammography to thyroid ultrasound,
and adds a triage layer on top. For each image:

\begin{enumerate}
  \item radiomics features are extracted from the image inside the nodule mask
        (Section~\ref{sec:methods-extraction});
  \item the features are standardised and a small subset is selected, using the training split
        only (Section~\ref{sec:methods-selection});
  \item the selected feature names and values are written as a short English text
        (Section~\ref{sec:methods-prompts});
  \item a pre-trained language model, fine-tuned for this task, reads the text and outputs a
        probability of malignancy (Section~\ref{sec:methods-llm});
  \item a three-way rule turns this score into \emph{auto-benign}, \emph{auto-malignant} or
        \emph{refer to a clinician}, with thresholds calibrated so that the error rates of the
        automatic decisions are statistically guaranteed (Section~\ref{sec:methods-triage}).
\end{enumerate}

An XGBoost classifier trained on the same features serves as the benchmark
(Section~\ref{sec:methods-xgb}). All code is in the thesis repository: the triage layer and
its unit tests in \texttt{conformal\_triage/} and \texttt{tests/}, the experiments in
\texttt{experiments/}, and one record per experiment in \texttt{runs/}.

\section{Radiomics feature extraction}
\label{sec:methods-extraction}

Features are extracted with PyRadiomics~\citep{vangriethuysen2017} in 2D mode, from the original
image only (no wavelet or Laplacian-of-Gaussian filtered images, as in~\citep{ra2025}). Each image
is converted to greyscale and rescaled to the range 0--255 by its own minimum and maximum, and
each mask is binarised at 127 (Section~\ref{sec:data-tn3k}). Because no pixel spacing is
available, the spacing is set to 1, so all sizes are in pixels.

\subsection{First extraction (v1)}

The first extraction used four feature families: shape (shape2D), first-order intensity
statistics, the grey-level co-occurrence matrix (GLCM) and the grey-level size-zone matrix
(GLSZM), giving 67 usable features. This is the set used by Ra et al. Texture features were
computed with a fixed bin width of 25 grey values.

The quality checks of Section~\ref{sec:data-qc} revealed three problems with this extraction:

\begin{itemize}
  \item with a bin width of 25, a nodule contains a median of about 9 grey levels (interquartile
        range 7--10), which leaves little room for texture;
  \item shape features in pixels partly measure the size of the image;
  \item multi-component masks were used as they are, so a single shape feature could describe
        several separate regions.
\end{itemize}

\subsection{Revised extraction (v2)}

The revised extraction changes four settings and keeps everything else:

\begin{itemize}
  \item \textbf{Fixed bin count.} Texture features use 64 grey levels per image instead of a fixed
        bin width.
  \item \textbf{Size relative to the image.} With image width $W$ and height $H$, lengths
        (major and minor axis, maximum diameter, perimeter) are divided by $\sqrt{WH}$, areas by
        $WH$, and the perimeter-to-surface ratio is multiplied by $\sqrt{WH}$. Shape features
        without units (sphericity, elongation) are unchanged.
  \item \textbf{Three more texture families:} the grey-level run-length matrix (GLRLM), the
        grey-level dependence matrix (GLDM) and the neighbouring grey-tone difference matrix
        (NGTDM).
  \item \textbf{Largest component.} For multi-component masks only the largest region is kept,
        a rule fixed before extraction so that each image describes one region.
\end{itemize}

This gives 102 features for all 3,493 images, with no extraction errors. As a consistency check,
on single-component masks the features shared with v1 reproduce the v1 values exactly.

\section{Preprocessing and feature selection}
\label{sec:methods-selection}

All steps below are fitted on the training split and then applied unchanged to the other splits:

\begin{enumerate}
  \item infinite values are set to missing and constant features are removed;
  \item missing values are replaced by the training median, and every feature is standardised to
        zero mean and unit variance (z-score);
  \item \textbf{near-duplicate filter} (later experiments): of two features with absolute
        Pearson correlation above 0.95, only the one more correlated with the label is kept;
  \item \textbf{mRMR}~\citep{peng2005} (minimum redundancy, maximum relevance) ranks the remaining
        features, and the first $k$ are used: $k = 8$ in the first experiments, as in Ra et al.,
        and $k = 16$ later.
\end{enumerate}

The near-duplicate filter was added after mRMR selected two features that are the same number in
2D (pixel surface and mesh surface, both relative to the image). The mRMR implementation penalises
a candidate by its \emph{average} correlation with the features already chosen, so a copy of one
chosen feature can still be picked when it is uncorrelated with the others. The selected features
keep their mRMR order in the prompt.

\section{From features to prompts}
\label{sec:methods-prompts}

Each nodule becomes one text: a fixed instruction followed by one sentence per selected feature.
Feature names are made readable (\texttt{original\_shape2D\_Elongation} becomes
\emph{Elongation}). With the format used in all main experiments (\emph{decimal}), a prompt begins:

\begin{quote}\small
Classify this thyroid nodule as benign or malignant based on the following standardized radiomic
features. The 'Major Axis Length Relative' feature is measured at 0.335. The 'Elongation' feature
is measured at $-$1.140. \dots
\end{quote}

The value is the z-score from Section~\ref{sec:methods-selection}, written with three decimals.
Four formats were compared (Table~\ref{tab:methods-formats}); a fifth is implemented but has not
been run yet (Section~\ref{sec:future-running}).

\begin{table}[ht]
  \centering
  \caption{Value formats in the prompt.}
  \label{tab:methods-formats}
  \begin{tabularx}{\linewidth}{l>{\raggedright\arraybackslash}X}
    \toprule
    format & sentence for one feature \\
    \midrule
    decimal    & The '\dots' feature is measured at $-$1.140. \\
    z-score    & The '\dots' feature has a z-score of $-$1.14 relative to the training nodules. \\
    percentile & The '\dots' feature is at percentile 12 of the training nodules. \\
    \midrule
    rich & z-score, percentile and a plain-language band, e.g.\ ``lower than
                         typical, meaning a less round, more irregular outline'' \\
    semantic only (not run yet) & the band and phrase without numbers \\
    \bottomrule
  \end{tabularx}
\end{table}

\emph{Decimal} and \emph{z-score} carry the same number; they differ only in saying what the
number is. \emph{Percentile} replaces the z-score by its rank among the training nodules, which
removes the information on how extreme a value is. The phrases of the two planned formats
describe what a feature measures, never which direction is malignant.

\section{Language-model classifier}
\label{sec:methods-llm}

\paragraph{Architecture.} The language model is used as a classifier, not as a text generator. A
linear classification head with two outputs is placed on the hidden state of the last token and
initialised at random; the probability of malignancy is the softmax of its two outputs. Ra et al.\
also train a classification head, with a binary cross-entropy loss, and compare it with training
the head alone.

\paragraph{Fine-tuning.} The pre-trained weights are loaded in 4-bit precision (NF4 with double
quantisation) and kept frozen; only low-rank adapters (LoRA) and the classification head are
trained (QLoRA~\citep{hu2022lora,dettmers2023qlora}). Table~\ref{tab:methods-llm} lists the
settings.

\begin{table}[ht]
  \centering
  \small
  \caption{Fine-tuning settings, with those reported by Ra et al.~\citep{ra2025}.}
  \label{tab:methods-llm}
  \begin{tabular}{lp{6.2cm}p{3.6cm}}
    \toprule
    setting & this work & Ra et al. \\
    \midrule
    backbone        & Qwen2.5-1.5B~\citep{qwen2024}; also Qwen2.5-7B and Llama-2-7B & Llama-2-7B~\citep{touvron2023llama2} \\
    weights         & 4-bit, frozen & not stated \\
    LoRA            & rank 16, $\alpha = 32$, dropout 0.05 & rank 16 \\
    LoRA placement  & attention Q and V (all linear layers in one run) & all layers \\
    loss            & class-weighted cross-entropy & binary cross-entropy \\
    optimiser       & AdamW, learning rate $2 \times 10^{-4}$, cosine schedule, 5\% warm-up & AdamW, $10^{-4}$, cosine schedule \\
    batch           & 8 $\times$ 2 accumulation steps (7B: 4 $\times$ 4) & not stated \\
    epochs          & 5 & over 500 \\
    model selection & epoch with the best validation AUC & best on validation \\
    \bottomrule
  \end{tabular}
\end{table}

Class weights are $n / (2 n_c)$ for a class with $n_c$ of the $n$ training nodules, so that
errors on the less frequent malignant class cost more. The maximum prompt length is checked
before training so that no prompt is truncated.

\paragraph{Hardware and precision.} The first models were trained on a free Colab T4 GPU in
16-bit floating point (fp16), the later ones on an L4 GPU in bfloat16. Probabilities are computed
in full precision after training. Earlier runs predicted in bfloat16, which rounded the scores so
that many nodules shared the same value (195 distinct values among 524 calibration nodules); this
barely changes the AUC but coarsens the thresholds of the triage layer.

\section{Classical benchmark}
\label{sec:methods-xgb}

XGBoost~\citep{chen2016xgboost} is trained on the same split, with the same preprocessing and the
same mRMR order, using 8, 16, 32 or all features. Its settings were fixed and not tuned: 300 trees
of depth 3, learning rate 0.05, and class weighting. The first experiments also used logistic
regression, a random forest and a support vector machine. XGBoost is a benchmark only; the
language model remains the model of interest.

\section{Conformal triage layer}
\label{sec:methods-triage}

\subsection{Three-way decision rule}

Let $s(x)$ be a score where higher means more likely malignant. With two thresholds $t < u$:
\[
  \text{decision}(x) =
  \begin{cases}
    \text{auto-benign}    & s(x) \le t, \\
    \text{auto-malignant} & s(x) \ge u, \\
    \text{refer}          & t < s(x) < u.
  \end{cases}
\]
A referred nodule follows the usual clinical path (radiologist review or biopsy). The model is not
changed by this layer; only $t$ and $u$ are chosen, on the calibration split.

\subsection{Risks and Learn-then-Test calibration}

Two clinical risks are controlled:

\begin{itemize}
  \item $R_1$: the share of malignant nodules among the auto-benign ones is at most
        $\alpha_1 = 5\%$ (cancers wrongly cleared);
  \item $R_2$: the share of benign nodules among the auto-malignant ones is at most
        $\alpha_2 = 20\%$ (unneeded biopsies or treatment).
\end{itemize}

The thresholds are chosen with Learn then Test (LTT)~\citep{angelopoulos2021ltt}, which turns
threshold selection into a sequence of hypothesis tests. For a candidate $t$, let $n$ be the
number of calibration nodules with $s \le t$ and $k$ the number of malignant ones among them. The
null hypothesis ``the true risk exceeds $\alpha_1$'' is tested with the binomial p-value
\[
  p(t) = P\big(\mathrm{Bin}(n, \alpha_1) \le k\big).
\]
Candidates are fixed in advance and tested from the safest to the most permissive. Testing stops
at the first $p(t) > \delta$, and the last certified candidate is kept (fixed-sequence testing,
which needs no multiplicity correction). $u$ is chosen the same way with $\alpha_2$. With
$\delta = 0.05$ for each risk, both guarantees hold together with probability at least
$1 - 2\delta = 0.90$ over the draw of the calibration set, provided calibration and future nodules
are exchangeable (Section~\ref{sec:data-official}). If no candidate is certified, that side of the
rule makes no automatic decisions.

The test requires a group that is both large and clean. With $\alpha_1 = \delta = 0.05$, an
auto-benign group needs at least 59 nodules with no cancer, 93 nodules with one cancer and 153
nodules with three cancers to be certified.

For the softmax probability the candidates are $t \in \{0.05, 0.10, \dots, 0.45\}$ and
$u \in \{0.95, 0.90, \dots, 0.55\}$.

\subsection{Uncertainty scores}
\label{sec:methods-scores}

\paragraph{Softmax probability.} The model's probability of malignancy for the unchanged prompt.

\paragraph{Perturbation votes.} Following the consistency-based uncertainty literature for
language models~\citep{ashby2026,noorani2025}, the original design of this work used an
agreement signal: the same nodule is presented to the unchanged model $m$ times in slightly
different ways, and the score is the number of queries $v$ answered malignant (probability
$\ge 0.5$). The first query is always the exact training prompt. The candidates for $m = 10$ are
$t \in \{0, \dots, 4\}$ and $u \in \{10, \dots, 6\}$. Two kinds of perturbations were used:

\begin{itemize}
  \item \textbf{Run 1 (mixed):} each query shuffled the feature order, dropped one feature with
        probability 0.5 and used one of three sentence templates.
  \item \textbf{Run 2 (format-preserving):} after run 1 showed that the model reacts to the format
        rather than to the content (Section~\ref{sec:results-ltt}), the template and order were
        kept. \emph{Jitter votes} add Gaussian noise to every standardised value (9 queries per
        noise level, standard deviation 0.05, 0.1 or 0.2, the level chosen on validation by a
        rule fixed in advance). \emph{Drop votes} remove one feature at a time (8 queries).
\end{itemize}

\subsection{Comparison with split conformal prediction}
\label{sec:methods-splitcp}

The same scores were also used with classical split conformal prediction~\citep{vovk2005}. The
nonconformity of a label is one minus its predicted probability, and a label enters the prediction
set when its nonconformity is below a threshold calibrated so that the true label is in the set
for at least $1-\alpha$ of the nodules. A set with one label counts as an automatic decision; a set
with both labels sends the nodule to a clinician. In the \emph{marginal} form one threshold serves
both labels; in the class-conditional (\emph{Mondrian}) form each label is calibrated on its own
nodules, so the guarantee holds for benign and malignant nodules separately. On the official split
both forms were run at $\alpha = 0.10$. On the pooled split the marginal form was run at
$\alpha = 0.10$ and $0.05$, and the Mondrian form with $\alpha = 0.10$ for benign and $0.05$ for
malignant nodules: then at most about 5\% of cancers receive the single label ``benign'', a
guarantee on missed cancers like the miss-rate variant below. Neither form bounds the share of
cancers \emph{among} the single-label benign sets, which is what $R_1$ requires.

\subsection{Exploratory miss-rate variant}
\label{sec:methods-miss}

$R_1$ is a statement about each auto-benign decision. A weaker, programme-level alternative bounds
the \emph{miss rate}: the share of all cancers sent to auto-benign, as in conformal risk
control~\citep{angelopoulos2024crc}. The same binomial test is used with $n$ equal to the number of
calibration cancers, on the grid $t \in \{0.01, 0.02, \dots, 0.50\}$, at 5\% and 10\%. This
variant was added after the main results and is reported as exploratory.

\section{Evaluation protocol}
\label{sec:methods-eval}

\begin{itemize}
  \item \textbf{Roles of the splits.} Training fits the model and all preprocessing; validation
        selects the epoch and the noise level; calibration is used only for the conformal
        thresholds; test only for reporting, with the rules locked.
  \item \textbf{Classification.} The main measure is the area under the ROC curve (AUC), computed
        from the predicted probability. Accuracy and sensitivity use the 0.5 threshold.
  \item \textbf{Comparing two runs.} Differences in AUC are given with a 95\% paired bootstrap
        interval: the same nodules are resampled 2,000 times and the AUC difference of the two
        runs is recomputed on each resample.
  \item \textbf{Separating errors.} How well a confidence score separates correct from incorrect
        predictions is measured by the AUC of that score for the event ``prediction correct''.
  \item \textbf{Exchangeability.} The classifier two-sample test of Section~\ref{sec:data-official}.
\end{itemize}

Unless stated otherwise, every run uses seed 42 and a single training run per setting.
